\documentclass[11pt]{amsart}

\usepackage[margin=1in]{geometry}
\usepackage{amsmath,amssymb,amsthm,mathtools}
\usepackage{microtype}
\usepackage{booktabs}
\usepackage{enumitem}
\usepackage{xcolor}
\usepackage{hyperref}
\usepackage[nameinlink,capitalize,noabbrev]{cleveref}

\newtheorem{theorem}{Theorem}[section]
\newtheorem{proposition}[theorem]{Proposition}
\newtheorem{corollary}[theorem]{Corollary}
\newtheorem{lemma}[theorem]{Lemma}
\theoremstyle{definition}
\newtheorem{definition}[theorem]{Definition}
\newtheorem{remark}[theorem]{Remark}

\newcommand{\cnum}{c}
\newcommand{\dist}{\operatorname{dist}}
\newcommand{\diam}{\operatorname{diam}}
\newcommand{\eps}{\varepsilon}
\newcommand{\whp}{\text{with high probability}}
\newcommand{\Occ}{\operatorname{Occ}}
\newcommand{\iotae}{\iota}

\hypersetup{
  colorlinks=true,
  linkcolor=blue!50!black,
  citecolor=blue!50!black,
  urlcolor=blue!50!black,
  pdftitle={Ball-Occupation Certificates under Coarse Graph Projections},
  pdfauthor={Levi Neuwirth},
  pdfkeywords={Cops and Robber, Meyniel's conjecture, graph products, degree reduction, expansion}
}

\title[Ball-Occupation Certificates under Coarse Graph Projections]{Ball-Occupation Certificates under Coarse Graph Projections\\
Degree Reduction, Square-Root Hard Families, and Toroidal Barriers}
\author{Levi Neuwirth}
\address{Brown University}
\date{July 2026}
\subjclass[2020]{05C57, 05C40, 05C12}
\keywords{Cops and Robber; Meyniel's conjecture; graph products; degree reduction; coarse graph projections; expansion; occupation certificates}

\begin{document}

\begin{abstract}
We isolate an abstract strategy-transfer principle for Cops and Robber.  Let
$\pi:V(H)\to V(G)$ have fibers of order at most $P$, and suppose that for a
scale factor $\lambda\ge1$ the distances between distinct fibers satisfy
\[
 \lambda(\dist_G(u,v)-1)+1
 \le \dist_H(x,y)
 \le \lambda(\dist_G(u,v)+2)-2,
\]
while every fiber has diameter at most $2(\lambda-1)$.  If the $N$-vertex
base graph has uniform multi-source growth at the square-root scale, then a
Hall assignment occupies a lifted macro-ball before the robber can leave it,
giving
\[
 \cnum(H)\le C P\bigl(d^3+\log(ePN)\bigr)\sqrt N
\]
and capture time at most $\lambda R$.  The theorem uses only the displayed
projection properties, not the internal form of the fibers.

Iterated degree reduction of Hosseini--Mohar--Gonzalez Hermosillo de la Maza
has exactly this geometry, with $\lambda=3^k$.  Applied to
$G(N,p)$ of expected degree $(\log N)^4$, it yields connected subcubic graphs
of order $M$ with
\[
 h(H)\ge (\log M)^{-O(1)},
 \qquad
 M^{\frac12-\frac{9+o(1)}{2\log\log\log M}}
 \le \cnum(H)
 \le \sqrt M\,(\log M)^{O(1)}.
\]
Thus the known stressing family has the square-root exponent, with a
polylogarithmically tight upper bound; the lower convergence to $1/2$ is only
triple-logarithmic.

Finally, we prove a sharp limitation of the strategy class.  If a common
prepositioned bank must, after learning a robber start $v$, match distinct cops
to every vertex of $B(v,R)$ using travel distance at most $R$, then
\[
 \Occ_R(G)
 \ge
 \frac{\sum_v |B(v,R)|}{\max_x |B(x,2R)|}.
\]
For every fixed $k\ge2$, the Cartesian torus $C_L^{\square k}$ has
$|V|=L^k$, vertex expansion $\Theta_k(L^{-1})$, exact cop number $k+1$, and
$\Occ_R=\Omega_k(L^k)$ for every radius $R$.  Hence, for every $\delta>0$,
there are bounded-degree graphs with $h(G)\ge |G|^{-\delta}$ and constant cop
number on which one-shot occupation still costs a linear number of cops.  A
cubic four-cycle replacement gives a degree-three instance at exponent
$\delta=1/2$.  Any universal robustness theorem must therefore use adaptive
multi-round pursuit, cop reuse, or another strategy not reducible to one-shot
occupation.
\end{abstract}

\maketitle

\section{Introduction and main conclusions}

The multi-cop version of Cops and Robber was developed by Aigner and Fromme,
who proved that three cops suffice on every planar graph
\cite{AignerFromme}.  Meyniel's conjecture asks whether every connected
$n$-vertex graph has cop number $O(\sqrt n)$.  The current best universal
upper bound remains
\[
 \frac{n}{2^{(1-o(1))\sqrt{\log_2 n}}},
\]
proved independently by Lu--Peng and Scott--Sudakov
\cite{LuPeng,ScottSudakov}.  Bose--Esperet--Hodor--Joret--Micek--Rambaud
recently extended the same scale of bound from graph order to vertex-cover
number \cite{CurrentFrontier}.  Expansion is one of the principal settings in
which polynomial savings are known: Bradshaw--Hosseini--Mohar--Stacho obtain
weak Meyniel bounds from bounded-degree expansion restricted to sublinear set
scales \cite{BHMS}, while Clow's withdrawn preprint developed a closely
related structural program connecting failure of weak Meyniel to high-cop
expanding examples \cite{Clow}.

The motivation here is the effect of bounded-degree replacement gadgets on
pursuit.  The degree-reduction construction of Hosseini--Mohar--Gonzalez
Hermosillo de la Maza (HMGHM) preserves lower bounds on cop number and
produces subcubic graphs with cop number $M^{1/2-o(1)}$ \cite{HMGHM}.  A
natural converse question is whether a useful upper strategy on the base
graph survives the replacement tower.

An arbitrary winning strategy does not lift transparently.  Moving one step
in the quotient may require a squad dispersed through a cloud to reorganize
while the robber continues moving.  The successful object is narrower and
more stable: an \emph{occupation certificate} that assigns distinct cops to
all vertices of a region before the robber can leave it.  Distance stretching
slows both deployment and escape, and a bounded normalized additive error
leaves a strict timing margin.

The first result is therefore stated for an abstract projection, rather than
for the HMGHM gadget.  The gadget enters only later, through an exact metric
calculation.  The resulting upper bound is stronger quantitatively than the
notation $M^{1/2+o(1)}$ suggests: it is $\sqrt M$ times a polylogarithmic
factor.  By contrast, the available lower bound approaches the square-root
exponent at the triple-logarithmic rate displayed in the abstract.  These two
facts should not be conflated merely because both can be written
$M^{1/2+o(1)}$.

The final result marks the boundary of the mechanism.  A one-shot occupation
certificate needs polynomial ball amplification between radii $R$ and $2R$.
Polynomially weak expansion alone does not supply this.  For every fixed
$k$, the Cartesian tori $C_L^{\square k}$ have bounded metric doubling, exact
cop number $k+1$, and linear one-shot occupation cost at every radius.  Taking
$k>1/\delta$ puts these examples inside every window
$h(G)\ge |G|^{-\delta}$.  A separate cubic replacement retains the barrier at
$\delta=1/2$.  The obstacle in the universal problem is therefore not degree
reduction itself; it is the need for adaptive reuse over many weak-growth
layers.

\section{Scale-adaptive cores and the robustness window}

For a connected graph $J$, write
\[
 h(J)=\min_{\varnothing\ne A\subseteq V(J),\ |A|\le |J|/2}
 \frac{|\partial_J A|}{|A|}.
\]
The following elementary reduction explains why polynomially weak expansion
is the relevant robustness window for the universal problem.

\begin{proposition}[Scale-adaptive induced core]\label{prop:adaptive-core}
Let $J$ be a connected graph of order $n$, and let
$\eta(1)\ge\cdots\ge\eta(n)\ge0$.  Then $J$ contains a connected induced
subgraph $K$, of order $m$, such that
\[
 \boxed{h(K)\ge\eta(m)}
 \qquad\text{and}\qquad
 \boxed{\cnum(J)\le \cnum(K)+\sum_{j=m+1}^{n}\eta(j).}
\]
\end{proposition}

\begin{proof}
Maintain the connected induced region $J_i$ containing the robber.  Whenever
$|J_i|=m_i$ and $h(J_i)<\eta(m_i)$, choose
$A_i\subseteq V(J_i)$ with $0<|A_i|\le m_i/2$ and
$|\partial_{J_i}A_i|<\eta(m_i)|A_i|$, and occupy its boundary.  The robber is
then confined to one component $J_{i+1}$ of
$J_i-\partial_{J_i}A_i$.  Whether that component lies inside $A_i$ or outside
it, one has
\[
 |A_i|\le m_i-m_{i+1}.
\]
Consequently the separator costs at most
\[
 \eta(m_i)(m_i-m_{i+1})
 \le
 \sum_{j=m_{i+1}+1}^{m_i}\eta(j).
\]
These integer intervals are disjoint along the robber's nested component
chain.  The process terminates at a connected induced $K$ with
$h(K)\ge\eta(|K|)$, and the separator costs telescope to the displayed sum.
\end{proof}

Taking $\eta(j)=j^{-a}$ shows that a polynomial cop-number saving on
subcubic graphs with $h(K)\ge |K|^{-a}$ would imply a weak form of Meyniel for
arbitrary graphs after the bounded-degree transfer of
Hosseini--Mohar--Gonzalez Hermosillo de la Maza
\cite[Corollary~8]{HMGHM}.  Bradshaw--Hosseini--Mohar--Stacho already treat
constant expansion restricted to sublinear set scales \cite{BHMS}; the
unresolved axis in this reduction is expansion that itself shrinks
polynomially.

\section{Coarse occupation projections}

\begin{definition}[Coarse occupation projection]\label{def:projection}
Let $G$ and $H$ be connected graphs.  A surjection
$\pi:V(H)\to V(G)$ is a $(\lambda,P)$-occupation projection if, writing
$F_v=\pi^{-1}(v)$,
\begin{enumerate}[label=\textup{(\roman*)}]
\item $|F_v|\le P$ for every $v\in V(G)$;
\item for distinct $u,v\in V(G)$ and arbitrary $x\in F_u$, $y\in F_v$,
\[
 \lambda(\dist_G(u,v)-1)+1
 \le
 \dist_H(x,y)
 \le
 \lambda(\dist_G(u,v)+2)-2;
\]
\item $\diam_H(F_v)\le2(\lambda-1)$ for every $v\in V(G)$.
\end{enumerate}
\end{definition}

The particular constants in Definition~\ref{def:projection} are chosen because they are
exact for the HMGHM tower.  The proof below only needs a bounded additive
slack after division by $\lambda$ and a strict gap between deployment and
escape deadlines.

For $U\subseteq V(G)$, let $B_G(U,r)$ be its closed radius-$r$
neighborhood.

\begin{theorem}[Abstract macro-ball occupation transfer]\label{thm:abstract-transfer}
Let $G$ have $N$ vertices.  Fix $d\ge2$, $R\ge2$, and constants $a,A_0>0$.
Assume
\begin{equation}\label{eq:scale}
 \sqrt N\le d^{R-2}<d\sqrt N,
 \qquad
 d^3\le\sqrt N,
\end{equation}
and
\begin{align}
 |B_G(U,R-2)|
 &\ge a\min\{|U|d^{R-2},N\}
 &&\text{for every }U\subseteq V(G),\label{eq:lowergrowth}\\
 |B_G(v,R)|
 &\le A_0d^R
 &&\text{for every }v\in V(G).\label{eq:uppergrowth}
\end{align}
If $H$ admits a $(\lambda,P)$-occupation projection onto $G$, then
\[
 \boxed{
 \cnum(H)
 \le
 C(a,A_0)P\bigl(d^3+\log(ePN)\bigr)\sqrt N.
 }
\]
The displayed number of cops captures the robber within at most $\lambda R$
cop moves.  In particular, the cop-number bound is independent of the scale
factor $\lambda$; tower depth affects capture time but not the required bank
size.
\end{theorem}

\begin{proof}
Put
\[
 \Theta=d^3+\log(ePN),
 \qquad
 \mu=\frac{AP\Theta}{\sqrt N},
\]
where $A$ is a sufficiently large constant depending only on $a,A_0$.
Choose one canonical vertex $z^*\in F_z$ in every fiber.  At $z^*$ place an
independent Poisson number of cop tokens of mean $\mu$.

For a possible robber fiber $F_v$, set
\[
 X_v=\pi^{-1}(B_G(v,R)).
\]
By \eqref{eq:uppergrowth} and \eqref{eq:scale},
\begin{equation}\label{eq:Qbound}
 Q_v:=|X_v|
 \le PA_0d^R
 <A_0Pd^3\sqrt N
 \le A_0P\Theta\sqrt N.
\end{equation}
We prove simultaneously for every $v$ that the sampled tokens can be matched
distinctly to all vertices of $X_v$, with every assigned token based over a
base vertex within distance $R-2$ of its target fiber.

Let $S\subseteq X_v$, $|S|=s$, and put $U=\pi(S)$.  Since every fiber has at
most $P$ targets, $|U|\ge s/P$.  Every token based over
$Z=B_G(U,R-2)$ is adjacent in the assignment graph to at least one target in
$S$.

If $|U|d^{R-2}<N$, then the number of available tokens is Poisson with mean
at least
\[
 \mu a|U|d^{R-2}\ge Aa\Theta s.
\]
For a Poisson variable $Y$ of mean $\Lambda\ge Aa\Theta s$,
\[
 \Pr(Y<s)\le e^{-\Lambda}\left(\frac{e\Lambda}{s}\right)^s.
\]
After increasing $A$, this is at most $(ePN)^{-4s}$.  There are at most
$N\binom{PN}{s}$ choices of a root and a target subset of size $s$, so summing
over $s\ge1$ gives $o(1)$.

If $|U|d^{R-2}\ge N$, then the available-token mean is at least
\[
 \mu aN=AaP\Theta\sqrt N
 \ge\frac{Aa}{A_0}Q_v
\]
by \eqref{eq:Qbound}.  Taking $A$ large and applying the same Poisson bound
shows, after a union bound over at most $N2^{Q_v}$ target subsets, that no
saturated Hall condition fails with probability $1-o(1)$.  Here every
saturated target set has
\[
 s\ge |U|\ge\frac{N}{d^{R-2}}>\frac{\sqrt N}{d}\ge N^{1/3},
\]
so the polynomial prefactors are negligible.

Hall's condition therefore holds simultaneously for all macro-balls with
probability $1-o(1)$.  The total number of sampled tokens is at most
$2AP\Theta\sqrt N$ with probability $1-o(1)$, so a deterministic placement of
the claimed size exists.

It remains to compare deadlines.  A target $y\in X_v$ lies over some
$u\in B_G(v,R)$.  Its assigned cop begins over $z$ with
$\dist_G(z,u)\le R-2$.  If $z\ne u$, the upper distortion bound gives travel
time at most
\[
 \lambda((R-2)+2)-2=\lambda R-2.
\]
If $z=u$, the fiber-diameter bound gives at most
$2(\lambda-1)\le\lambda R-2$, since $R\ge2$.

To leave $X_v$, the robber must enter a fiber over a base vertex at distance
at least $R+1$ from $v$.  The lower distortion bound makes this require at
least
\[
 \lambda((R+1)-1)+1=\lambda R+1
\]
steps.  Every vertex of $X_v$ is occupied first, and the cop assigned to the
robber's current vertex captures her.
\end{proof}

\begin{remark}[The quantifier needed from the random base]\label{rem:PW-uniform}
The use of \eqref{eq:lowergrowth} is graph-uniform, not a per-source-set
probability statement.  In the dense theorem of Pra{\l}at and Wormald,
condition~(i) of their deterministic Theorem~3.1 is explicitly quantified over
\emph{every} source set and radius.  Their Theorem~3.4 proves that a single
$G(N,p)$ satisfies those hypotheses asymptotically almost surely; its proof
unions over the bad source sets and concludes that the growth estimate holds
simultaneously for all sets and radii \cite[Theorems~3.1 and~3.4]{PW}.
Thus the external input has the quantifier order required by
\cref{thm:abstract-transfer}.
\end{remark}

\section{The HMGHM replacement tower}

For a vertex of degree $r$, the HMGHM replacement has one external port for
every incident edge.  The ports are partitioned into nearly equal classes,
and for each pair of classes there is an internal vertex adjacent to every
port in the two classes \cite[Section~2]{HMGHM}.

\begin{lemma}[One-round port geometry]\label{lem:portgeometry}
For every HMGHM replacement cloud of degree at least two:
\begin{enumerate}[label=\textup{(\roman*)}]
\item distinct ports are nonadjacent and have distance exactly two;
\item every cloud vertex is within distance at most three of every specified
port;
\item the cloud diameter is at most four.
\end{enumerate}
\end{lemma}

\begin{proof}
Two ports in different classes share the internal vertex associated with
their class pair.  Two ports in the same class share any internal vertex
associated with that class and another nonempty class.  Since ports are
mutually nonadjacent, their distance is exactly two.

An internal vertex is adjacent to every port in either of two classes.  If a
specified port lies in neither class, travel to a port in one of the two
classes, then through the internal vertex corresponding to that class and the
specified port's class, and finally to the specified port.  This takes three
steps.  The diameter bound follows by routing arbitrary endpoints through a
specified port.
\end{proof}

Let
\[
 G=G_0,G_1,\ldots,G_k=H
\]
be an iterated HMGHM tower, and let $\pi:V(H)\to V(G)$ map every final vertex
to its original ancestor.  Put
\[
 F_v=\pi^{-1}(v),
 \qquad
 P=\max_v|F_v|,
 \qquad
 \lambda=3^k.
\]

\begin{theorem}[Exact normalized distortion]\label{thm:metric}
The ancestry projection is a $(3^k,P)$-occupation projection.  Explicitly,
for distinct base vertices $u,v$, $r=\dist_G(u,v)$, and arbitrary
$x\in F_u$, $y\in F_v$,
\[
 \boxed{
 3^k(r-1)+1
 \le
 \dist_H(x,y)
 \le
 3^k(r+2)-2,
 }
\]
and
\[
 \boxed{\diam_H(F_v)\le2(3^k-1).}
\]
\end{theorem}

\begin{proof}
For one round, a shortest path between distinct clouds uses $e\ge r$
external edges.  Between consecutive external edges it enters and leaves an
intermediate cloud through distinct ports: otherwise it immediately traverses
one external edge back.  By Lemma~\ref{lem:portgeometry}, each intermediate port
change costs at least two internal edges, and hence
\[
 \dist_{G_1}(x,y)\ge e+2(e-1)\ge3r-2.
\]
For the upper bound, follow a base geodesic.  Reaching the first prescribed
port costs at most three, each intermediate port change costs two, the
external edges cost $r$, and reaching the final endpoint costs at most three.
Thus
\[
 \dist_{G_1}(x,y)\le3+r+2(r-1)+3=3r+4.
\]
The one-round fiber diameter is at most four.

The lower and upper affine recurrences are
\[
 L_j(r)=3L_{j-1}(r)-2,
 \qquad
 U_j(r)=3U_{j-1}(r)+4,
\]
with $L_0(r)=U_0(r)=r$.  Solving gives
\[
 L_k(r)=3^k(r-1)+1,
 \qquad
 U_k(r)=3^k(r+2)-2.
\]
The diameter recurrence $D_j\le3D_{j-1}+4$, $D_0=0$, gives
$D_k\le2(3^k-1)$.
\end{proof}

The feature that matters is not the number of rounds but the normalized
additive error: after division by $3^k$, it remains two quotient layers.  By
\cref{thm:abstract-transfer}, any other graph projection with the same three
properties inherits the same occupation-certificate transfer.

\section{Expansion retention under port-cloud replacement}

\begin{proposition}[Expansion under connected port replacement]\label{prop:port-exp}
Let $H$ be obtained from a base graph $G$ by replacing every vertex by a
connected cloud of order at most $L_0$, with distinct external ports for the
incident base edges.  If $\iotae(G)$ is the edge-isoperimetric constant of
$G$, then
\[
 \boxed{
 \iotae(H)
 \ge
 \frac{1}{2L_0}
 \min\left\{1,\frac{\iotae(G)}{L_0}\right\}.
 }
\]
\end{proposition}

\begin{proof}
Let $S\subseteq V(H)$ with $0<|S|\le|H|/2$.  In each cloud, classify the
majority side and let $\mathcal M$ be the total number of minority vertices.  Since
every partially cut cloud is connected and has at most $L_0$ vertices, its
internal cut contributes at least one edge, so the total internal contribution
is at least $\mathcal M/L_0$.

Let $U$ be the set of base vertices whose clouds have majority in $S$, and
put $E=e_G(U,V(G)\setminus U)$.  A base cut edge can fail to cross the lifted
cut only if one of its two ports is a minority vertex.  Distinct base edges
use distinct ports, so at most $\mathcal M$ of the $E$ external cut edges fail.  Hence
\[
 e_H(S,V(H)\setminus S)
 \ge
 \frac{\mathcal M}{L_0}+\max\{0,E-\mathcal M\}
 \ge
 \frac{E+\mathcal M}{2L_0}.
\]
Apply the same majority accounting to $S$ or its complement, according as
$|U|\le|G|/2$ or not, to obtain
\[
 \min\{|U|,|V(G)\setminus U|\}
 \ge
 \frac{|S|-\mathcal M}{L_0}.
\]
Thus
\[
 E\ge\frac{\iotae(G)}{L_0}(|S|-\mathcal M),
\]
and consequently
\[
 E+\mathcal M
 \ge
 \min\left\{1,\frac{\iotae(G)}{L_0}\right\}|S|.
\]
Combining the inequalities proves the proposition.
\end{proof}

The HMGHM cloud-size recurrence turns the one-round estimate into a
subpolynomial-loss statement in the polylogarithmic-degree regime.  If $D$
is the initial maximum degree, $L_i$ is the largest cloud order in round $i$,
and $k$ rounds are used, HMGHM prove
\[
 \prod_{i=0}^{k-1}L_i
 \le C D^2(\log D)^{\log_2(11/5)},
 \qquad
 2^k=O(\log D).
\]
Iterating Proposition~\ref{prop:port-exp} therefore gives the following.

\begin{corollary}[Expansion retained by HMGHM reduction]\label{cor:exp-retention}
Let $H$ be the final subcubic graph obtained from a connected graph $G$ of
maximum degree $D\ge4$.  Then
\[
 \boxed{
 h(H)
 \ge
 \frac{\iotae(G)}{C D^4(\log D)^{\kappa}},
 \qquad
 \kappa=1+2\log_2(11/5)<3.28.
 }
\]
In particular, if $D=|G|^{o(1)}$ and $\iotae(G)=|G|^{-o(1)}$, then
$|H|=|G|^{1+o(1)}$ and $h(H)=|H|^{-o(1)}$.
\end{corollary}

\begin{proof}
At every round $\iotae(G_i)\le D_i\le L_i$, so
Proposition~\ref{prop:port-exp} gives
\[
 \iotae(G_{i+1})\ge\frac{\iotae(G_i)}{2L_i^2}.
\]
Thus
\[
 \iotae(H)
 \ge
 \frac{\iotae(G)}{2^k(\prod_iL_i)^2}
 \ge
 \frac{\iotae(G)}{C D^4(\log D)^\kappa}.
\]
Since $H$ has maximum degree at most three, its vertex expansion is at least
one third of its edge expansion.  The order statement follows from the same
cloud-product bound.
\end{proof}

\begin{remark}[Relation to replacement products]
The regular replacement-product literature proves stronger spectral
conclusions under much stronger hypotheses on the clouds; see, for example,
Reingold--Vadhan--Wigderson \cite{RVW}.  Proposition~\ref{prop:port-exp} allows arbitrary
connected, nonuniform clouds and consequently gives only a crude
isoperimetric estimate.  No novelty claim is made here beyond this precise
form without a fuller graph-substitution review.
\end{remark}

\section{A polylogarithmically tight square-root family}

Take
\[
 d=(\log N)^4,
 \qquad
 p=\frac{d}{N-1},
 \qquad
 G\sim G(N,p).
\]
Iterate the HMGHM replacement until the graph $H$ is subcubic, and write
$M=|H|$.

With high probability, $\Delta(G)\le2d$.  By
Remark~\ref{rem:PW-uniform}, the dense Pra{\l}at--Wormald theorem supplies the
uniform lower growth in \eqref{eq:lowergrowth}; in the volume range used here
it also supplies the upper growth in \eqref{eq:uppergrowth} \cite{PW}.  Choose
$R$ minimally so that $d^{R-2}\ge\sqrt N$.  Since $d$ is polylogarithmic,
\eqref{eq:scale} holds.

HMGHM give, both globally and along one ancestry fiber,
\begin{equation}\label{eq:Pbound}
 P\le C d^2(\log d)^{1.14},
 \qquad
 N\le M\le PN.
\end{equation}
Their shadow strategy gives $\cnum(H)\ge\cnum(G)$, and the random-graph lower
bound of Bollob\'as--Kun--Leader used in their argument yields
\[
 \cnum(G)
 \ge
 d^{-2}N^{\frac12-\frac{9}{2\log\log d}}.
\]
The abstract transfer theorem gives
\[
 \cnum(H)
 \le
 CP\bigl(d^3+\log(ePN)\bigr)\sqrt N
 \le
 \sqrt M\,(\log M)^{20+o(1)}.
\]
Using $d=(\log N)^4$ and $M=N(\log N)^{O(1)}$ in the lower bound gives the
following more informative formulation.

\begin{theorem}[Quantitative HMGHM hard family]\label{thm:hardfamily}
There is a sequence of connected subcubic graphs $H$, of order $M\to\infty$,
for which
\[
 \boxed{
 M^{\frac12-\frac{9+o(1)}{2\log\log\log M}}
 \le
 \cnum(H)
 \le
 \sqrt M\,(\log M)^{20+o(1)}.
 }
\]
The upper bound is $\sqrt M$ times a polylogarithmic factor.  The lower
exponent tends to $1/2$ only at a triple-logarithmic rate.
\end{theorem}

The base edge expansion is $\Omega(d)$ with high probability.  Since
$d=\operatorname{polylog}N$, Corollary~\ref{cor:exp-retention} gives
\[
 h(H)\ge (\log M)^{-O(1)}=M^{-o(1)}.
\]
This is exactly the degree regime in which the retention factor is
informative; for polynomial initial degree the crude $D^4$ loss can be
vacuous.

\begin{corollary}[Square-root weak-expander family]\label{cor:weakexpander}
There are connected subcubic graphs satisfying
\[
 h(H)\ge M^{-o(1)}
 \qquad\text{and}\qquad
 \cnum(H)=M^{1/2+o(1)}.
\]
More precisely, they obey the two-sided bounds of
\cref{thm:hardfamily}.
\end{corollary}

\begin{remark}[What is forced, and what is achieved]\label{rem:robustness-endpoint}
By \cref{prop:adaptive-core}, polynomially weak expansion is a natural
robustness window for weak Meyniel.  The HMGHM lower bound alone already
forces every proposed estimate
\[
 \cnum(J)\le C\phi^{-p}|J|^{1-\eps+o(1)}
 \qquad(h(J)\ge\phi=|J|^{-o(1)})
\]
to have $\eps\le1/2$.  That restriction predates the upper transfer proved
here.  The new conclusion is that the known stressing family itself achieves
the endpoint order $\sqrt M$ up to polylogarithmic factors, so this family is
not an obstruction to a Meyniel-strength theorem on polynomially weak
subcubic expanders.
\end{remark}

\section{Why chase strategies need not transfer}

The abstract theorem deliberately transfers a strategy class, not arbitrary
cop number.  The smallest example explains the distinction.  For a degree-two
vertex, one HMGHM cloud is a three-vertex path.  Replacing every vertex of
$C_3$ therefore produces $C_9$.  But
\[
 \cnum(C_3)=1,
 \qquad
 \cnum(C_9)=2.
\]
The one-cop win on $C_3$ is a direct chase/dismantling phenomenon.  The
subdivision-like stretching destroys it.  By contrast, an occupation
certificate is synchronized to a deadline: the replacement tower stretches
the cops' travel and the robber's escape by the same factor, and the bounded
normalized additive slack preserves a strict margin.  The examples
$K_4$ and the diamond graph exhibit the same one-round increase, so the issue
is structural rather than peculiar to one cycle.

\section{A one-shot occupation barrier}

\begin{definition}[Universal one-shot occupation number]\label{def:occ}
For a connected graph $G$ and integer $R\ge0$, let $\Occ_R(G)$ be the minimum
size of a finite set $X$ of distinct cop tokens, equipped with a position map
$p:X\to V(G)$, such that for every $v\in V(G)$ there is an injection
\[
 f_v:B_G(v,R)\longrightarrow X
\]
with
\[
 \dist_G(u,p(f_v(u)))\le R
 \qquad\text{for every }u\in B_G(v,R).
\]
Different tokens may have the same initial position.  After learning the
robber's starting vertex, the common prepositioned bank can occupy her entire
radius-$R$ ball within $R$ moves.
\end{definition}

\begin{remark}[Why the target and deadline are natural]
If the robber starts at $v$, she needs at least $R+1$ robber moves to leave
$B_G(v,R)$.  Occupying that whole ball within $R$ cop moves is therefore the
canonical one-shot certificate: every vertex she could still occupy is filled
before her first possible escape.  The parameter $\Occ_R$ measures this
specific strategy class, not ordinary cop number.
\end{remark}

\begin{theorem}[Counting barrier]\label{thm:occ-lower}
Every connected graph satisfies
\[
 \boxed{
 \Occ_R(G)
 \ge
 \frac{\sum_{v\in V(G)}|B_G(v,R)|}
 {\max_{x\in V(G)}|B_G(x,2R)|}.
 }
\]
In particular, if $G$ is vertex-transitive, then
\[
 \boxed{
 \Occ_R(G)
 \ge
 |V(G)|\frac{|B_G(o,R)|}{|B_G(o,2R)|}.
 }
\]
\end{theorem}

\begin{proof}
Fix a feasible multiset $X$.  For each possible robber start $v$, every cop
token used by the injection $f_v$ lies in $B_G(v,2R)$, by the triangle
inequality.  Hence at least $|B_G(v,R)|$ tokens of $X$ lie in $B_G(v,2R)$.
Summing over $v$, the number of incident pairs $(v,x)$ with
$x\in X\cap B_G(v,2R)$ is at least
$\sum_v|B_G(v,R)|$.

A fixed token based at $x$ is counted only for starts
$v\in B_G(x,2R)$, at most $\max_y|B_G(y,2R)|$ times.  Therefore
\[
 |X|\max_y|B_G(y,2R)|
 \ge
 \sum_v|B_G(v,R)|,
\]
which proves the claim.
\end{proof}

The theorem identifies the exact growth ratio demanded by one-shot
occupation.  A polynomial saving from the trivial $|V(G)|$ bound requires
polynomial amplification from radius $R$ to radius $2R$.

We now give subcubic witnesses showing that polynomially weak expansion does
not imply such amplification.

\begin{definition}[The cubic truncated torus]\label{def:Qt}
For $L\ge5$, let $Q_L$ have vertex set
\[
 (\mathbb Z/L\mathbb Z)^2\times\mathbb Z/4\mathbb Z.
\]
Inside each fiber $(x,y)\times\mathbb Z/4\mathbb Z$, join the four vertices in
a cycle.  Add the external edges
\[
 (x,y,0)(x,y+1,2)
 \qquad\text{and}\qquad
 (x,y,1)(x+1,y,3)
\]
for every $(x,y)$.  Equivalently, $(x,y,2)$ receives its external edge from
$(x,y-1,0)$, and $(x,y,3)$ receives its external edge from $(x-1,y,1)$.
Thus every vertex has two internal cycle neighbors and one external neighbor,
and $Q_L$ is the four-cycle port replacement of the square torus
$C_L\square C_L$.
\end{definition}

Every vertex of $Q_L$ has degree three, and the construction embeds on the
torus by replacing each base vertex inside a small disk.  Its order is
$4L^2$.



\begin{lemma}[Vertex transitivity and explicit doubling of $Q_L$]\label{lem:doubling}
The graph $Q_L$ is vertex-transitive and, for every vertex $x$ and radius
$R\ge0$,
\[
 |B_{Q_L}(x,2R)|\le 5500\,|B_{Q_L}(x,R)|.
\]
\end{lemma}

\begin{proof}
Translations in the first two coordinates are automorphisms.  The map
\[
 \rho(x,y,i)=(y,-x,i+1)
\]
(with coordinates interpreted cyclically) preserves internal cycle edges and
interchanges the two external edge directions.  Translations together with
$\rho$ act transitively.

Let $T_L=C_L\square C_L$ and project $(x,y,i)$ to $(x,y)$.  Projection does
not increase distance, so
\[
 |B_{Q_L}(x,2R)|\le4|B_{T_L}(\pi x,2R)|
 \le4\min\{L,4R+1\}^2.
\]
Conversely, from an arbitrary cloud vertex one can enter the required port in
at most two internal moves and then lift each base step using at most three
moves.  Hence, with $r=\lfloor(R-2)/3\rfloor$ for $R\ge2$,
\[
 |B_{Q_L}(x,R)|\ge |B_{T_L}(\pi x,r)|.
\]
The coordinate box of cyclic radius $\lfloor r/2\rfloor$ lies inside the
$\ell_1$ ball, so
\[
 |B_{T_L}(\pi x,r)|
 \ge \min\{L,2\lfloor r/2\rfloor+1\}^2.
\]
For $R<10$, the upper bound is at most $4\cdot37^2<5500$ and the denominator
is at least one.  For $R\ge10$, one has $r\ge R/6$ and
$2\lfloor r/2\rfloor+1\ge r$, whence the ratio is at most
$4\cdot30^2<5500$.  This proves the displayed constant.
\end{proof}

\begin{theorem}[Cubic one-shot barrier]\label{thm:cubic-barrier}
For $L\ge5$, the connected cubic graphs $Q_L$, with $M=|Q_L|=4L^2$, satisfy
\[
 \boxed{
 h(Q_L)=\Theta(M^{-1/2}),
 \qquad
 \cnum(Q_L)\le3,
 \qquad
 \Occ_R(Q_L)\ge \frac{M}{5500}
 \quad\text{for every }R\ge0.
 }
\]
\end{theorem}

\begin{proof}
The square torus has edge and vertex expansion $\Theta(1/L)$ by the discrete
torus isoperimetric inequality \cite{BL}.  Applying
Proposition~\ref{prop:port-exp} with cloud size four gives the matching lower
bound for $Q_L$; lifting a coordinate slab of width $\lfloor L/2\rfloor$
gives the upper bound for every $L$.  Since $Q_L$ is cubic, edge and vertex
expansion differ by at most a constant factor.
Thus $h(Q_L)=\Theta(1/L)=\Theta(M^{-1/2})$.

The graph $Q_L$ is toroidal, and every toroidal graph has cop number at most
three \cite{Lehner}.  Vertex transitivity, Theorem~\ref{thm:occ-lower}, and
Lemma~\ref{lem:doubling} give
\[
 \Occ_R(Q_L)
 \ge
 M\frac{|B_{Q_L}(x,R)|}{|B_{Q_L}(x,2R)|}
 \ge \frac{M}{5500}.
\]
The finite audit suggests that the optimal asymptotic constant is $1/4$, but
that sharpening is not needed here.
\end{proof}

\subsection{The barrier throughout every polynomial expansion window}

For integers $k\ge2$ and $L\ge4$, write
\[
 T_{L,k}=\underbrace{C_L\square\cdots\square C_L}_{k\text{ factors}}.
\]
Its order is $m=L^k$, its degree is $2k$, and its metric is the cyclic
$\ell_1$ metric.

\begin{lemma}[Uniform doubling of Cartesian tori]\label{lem:ktorus-doubling}
For every fixed $k\ge2$, every $L\ge4$, every vertex $x$, and every radius
$R\ge0$,
\[
 |B_{T_{L,k}}(x,2R)|\le (5k)^k|B_{T_{L,k}}(x,R)|.
\]
\end{lemma}

\begin{proof}
Every coordinate of a point in $B(x,2R)$ has cyclic distance at most $2R$, so
\[
 |B(x,2R)|\le \min\{L,4R+1\}^k.
\]
The coordinate box in which every coordinate has cyclic distance at most
$\lfloor R/k\rfloor$ lies in $B(x,R)$, and therefore
\[
 |B(x,R)|\ge\min\{L,2\lfloor R/k\rfloor+1\}^k.
\]
If $R<k$, the ratio is at most $(4k+1)^k$.  If $R\ge k$, then
$2\lfloor R/k\rfloor+1\ge R/k$ and $4R+1\le5R$; taking the minima with $L$
does not increase their ratio beyond $5k$.  The claim follows.
\end{proof}

\begin{theorem}[Full-window toroidal barrier]\label{thm:full-window}
For every fixed $k\ge2$, the graphs $T_{L,k}$ satisfy
\[
 \boxed{
 h(T_{L,k})=\Theta_k(m^{-1/k}),
 \qquad
 \cnum(T_{L,k})=k+1,
 \qquad
 \Occ_R(T_{L,k})\ge (5k)^{-k}m
 \quad\text{for every }R\ge0.
 }
\]
Consequently, for every $\delta>0$ there is a constant-degree graph family
with
\[
 h(G)\ge |G|^{-\delta},
 \qquad
 c(G)=O_\delta(1),
 \qquad
 \Occ_R(G)=\Omega_\delta(|G|)
 \quad\text{for every radius }R.
\]
\end{theorem}

\begin{proof}
The discrete-torus edge-isoperimetric inequality gives order $1/L$
\cite{BL}.  Since $T_{L,k}$ has degree $2k$, edge boundary and external
vertex boundary differ by at most the fixed factor $2k$; a coordinate slab of
width $\lfloor L/2\rfloor$ supplies the matching upper bound.  Hence
$h(T_{L,k})=\Theta_k(1/L)=\Theta_k(m^{-1/k})$.  Neufeld and Nowakowski proved
that a Cartesian product of $k$ cycles, each of length at least four, has cop
number exactly $k+1$ \cite{NeufeldNowakowski}.  Since the torus is
vertex-transitive, Theorem~\ref{thm:occ-lower} and
Lemma~\ref{lem:ktorus-doubling} give the occupation lower bound.

Given $\delta>0$, choose $k=\max\{2,\lfloor1/\delta\rfloor+1\}$.  Then $1/k<\delta$, so for sufficiently large $m$ the expansion
lower bound $h(T_{L,k})\ge m^{-\delta}$ holds after absorbing the fixed
$k$-dependent constant.
\end{proof}

\begin{remark}[The sharp metric constant]
For fixed $k$, choose radii $1\ll R\ll L$.  Lattice-point asymptotics for the
$\ell_1$ ball give
\[
 \frac{|B_{T_{L,k}}(x,2R)|}{|B_{T_{L,k}}(x,R)|}=2^k+o(1).
\]
Thus no uniform doubling constant below $2^k$ is possible, and the counting
bound of \cref{thm:occ-lower} approaches the natural fraction $2^{-k}m$ on
these local radii.  The explicit constant $(5k)^{-k}$ is chosen only for a
short all-radii proof.
\end{remark}

\begin{corollary}[No expansion-only one-shot theorem]\label{cor:no-exp-occ}
For every $\delta>0$, there is no implication of the form
\[
 h(G)\ge |G|^{-\delta}
 \quad\Longrightarrow\quad
 \Occ_R(G)\le |G|^{1-\eps}
 \text{ for some radius $R$}
\]
with any fixed $\eps>0$, even when the maximum degree is bounded by a constant
depending only on $\delta$.
\end{corollary}

\begin{proof}
Take the family from \cref{thm:full-window}.  It lies in the prescribed
expansion window, while $\Occ_R(G)=\Omega_\delta(|G|)$ for every $R$.
\end{proof}

\begin{remark}[Architectural meaning]
The logical obstruction is not that tori are difficult pursuit instances;
they are not.  Rather, \cref{thm:occ-lower} makes every vertex-transitive
bounded-doubling graph expensive for one-shot occupation, and bounded
doubling is compatible with every polynomial weak-expansion window by
\cref{thm:full-window}.  The tori
certify that the hypothesis class contains such graphs while coordinate-wise
shadowing still uses exactly $k+1$ cops.  Therefore ball amplification is the
wrong invariant for adaptive pursuit, not merely a poor description of one
particular easy family.  The cubic family $Q_L$ records that the same
separation already occurs in maximum degree three at exponent $1/2$.
\end{remark}

\section{Outlook}

For vertex expansion $h(G)\ge\phi$, iterating the elementary growth factor
$1+\phi$ reaches global scale after $O(\phi^{-1}\log |G|)$ layers.  The same
iteration gives the standard diameter bound of that order; these are two
forms of the same calculation, not independent evidence.  The strategic
consequence is that, when $\phi=|G|^{-a}$, an amplification-based pursuit
scheme must operate over a full-traversal timescale $\Theta(|G|^a\log|G|)$.

The present paper separates three phenomena:
\begin{enumerate}[label=\textup{(\arabic*)}]
\item bounded normalized metric distortion preserves a strong occupation
certificate through degree reduction;
\item the HMGHM stressing family itself meets the square-root endpoint up to
polylogarithmic factors;
\item one-shot occupation is nevertheless incapable of proving a universal
robustness theorem throughout any polynomial weak-expansion window, even on
bounded-degree graphs with constant cop number; a cubic instance already
appears at exponent $1/2$.
\end{enumerate}
The remaining universal question is therefore an adaptive one.  On the tori,
ball growth carries essentially no information about pursuit cost; product
structure instead supports coordinate-wise shadowing.  What geometric or
combinatorial quantity replaces product coordinates on a general
polynomially weak expander?  Equivalently, can a capacitated, correlated, or
deferred witness system reuse the same cop resources over polynomially many
weak-growth layers, or must every such one-traversal certificate incur
polynomial congestion?

\section*{Acknowledgments}
The author is grateful to Anthony Clow, Peter Bradshaw, Bojan Mohar, and
Florian Lehner for work and perspectives that helped shape the questions
addressed here.  Additional acknowledgments will be added in a later version.
The author welcomes corrections concerning priority, related
graph-substitution inequalities, and the scope of the occupation framework.

\section*{Audit and reproducibility}

The metric inequalities were independently tested on HMGHM towers rebuilt
from the published gadget description, including structured base graphs not
used in the original audit.  The Hall inequalities and timing margins were
checked numerically, and exact small replacement games were solved by
retrograde analysis.  The toroidal barrier audit computes exact ball profiles of
$C_L^{\square k}$ by convolving cyclic distance distributions, verifies the
$(5k)^k$ doubling bound for $k=2,3,4,5$, constructs $Q_L$, checks cubicity,
connectivity, and the displayed rotation automorphism, and evaluates the
counting lower bound at every radius.  No theorem depends on the
computations.

\begin{thebibliography}{99}

\bibitem{AignerFromme}
M.~Aigner and M.~Fromme,
\emph{A game of cops and robbers},
Discrete Applied Mathematics 8 (1984), 1--12.

\bibitem{CurrentFrontier}
P.~Bose, L.~Esperet, J.~Hodor, G.~Joret, P.~Micek, and C.~Rambaud,
\emph{Cops and robber in graphs with bounded vertex cover number},
arXiv:2602.07435, 2026.

\bibitem{BHMS}
P.~Bradshaw, S.~A. Hosseini, B.~Mohar, and L.~Stacho,
\emph{On the cop number of graphs of high girth},
Journal of Graph Theory 102 (2023), 15--34; arXiv:2005.10849.

\bibitem{BKL}
B.~Bollob\'as, G.~Kun, and I.~Leader,
\emph{Cops and robbers in a random graph},
Journal of Combinatorial Theory, Series B 103 (2013), 226--236.

\bibitem{BL}
B.~Bollob\'as and I.~Leader,
\emph{An isoperimetric inequality on the discrete torus},
SIAM Journal on Discrete Mathematics 3 (1990), 32--37.

\bibitem{Clow}
A.~Clow,
\emph{Expanders satisfy the weak Meyniel conjecture},
withdrawn preprint, arXiv:2311.13792, 2023.

\bibitem{HMGHM}
S.~A. Hosseini, B.~Mohar, and S.~Gonzalez Hermosillo de la Maza,
\emph{Meyniel's conjecture on graphs of bounded degree},
Journal of Graph Theory 97 (2021), 401--407; arXiv:1912.06957.

\bibitem{Lehner}
F.~Lehner,
\emph{On the cop number of toroidal graphs},
Journal of Combinatorial Theory, Series B 151 (2021), 250--262.

\bibitem{LuPeng}
L.~Lu and X.~Peng,
\emph{On Meyniel's conjecture of the cop number},
Journal of Graph Theory 71 (2012), 192--205.

\bibitem{NeufeldNowakowski}
S.~Neufeld and R.~Nowakowski,
\emph{A game of cops and robbers played on products of graphs},
Discrete Mathematics 186 (1998), 253--268.

\bibitem{PW}
P.~Pra{\l}at and N.~Wormald,
\emph{Meyniel's conjecture holds for random graphs},
Random Structures \& Algorithms 48 (2016), 396--421; arXiv:1301.2841.

\bibitem{RVW}
O.~Reingold, S.~Vadhan, and A.~Wigderson,
\emph{Entropy waves, the zig-zag graph product, and new constant-degree
expanders},
Annals of Mathematics 155 (2002), 157--187.

\bibitem{ScottSudakov}
A.~Scott and B.~Sudakov,
\emph{A bound for the cops and robbers problem},
SIAM Journal on Discrete Mathematics 25 (2011), 1438--1442.

\end{thebibliography}

\end{document}
